\section{Threats to Validity}\label{sec:threats}

\subsection{Construct Validity}

\paragraph{What a validity flag measures.}
The central construct threat in this area is that the usual ``valid suite'' indicator does
not measure what its name suggests, and fails silently when it does not. We encountered six
distinct instances, four of them in our own instrumentation.

First, a compilation flag inferred from a coverage report measures only whether the class
appears in that report. Coverage tools enumerate every class in a project, including classes
that were never loaded; such a class yields zero covered branches, which the parser then
records as a successfully compiled function with zero coverage. Second, a criterion of ``at
least one test executed'' counts a file in which every test fails. Third, reflection removes
the compiler's ability to check names, so a compilation-based rate is uninformative for any
suite that uses it. Fourth, mismatched tool versions can make a test launcher discover zero
tests while exiting successfully.

The fifth is worth stating in full, because believing it would have inverted a published
conclusion. Our Java runner and the generated tests were compiled with the default JDK
(class file version 61) and executed under the JDK~11 required by EvoSuite (which reads at
most 55). The resulting \texttt{UnsupportedClassVersionError} was raised inside a child
process, so the harness observed only the absence of a result line and recorded ``no tests
ran'' --- for every target. EvoSuite scored $0/60$, and a paired test against it returned
$p = 0.005$ with a rank-biserial correlation of $-1.000$. Had we reported it, the paper would
have claimed a statistically significant advantage for the model over EvoSuite; the true
figure is given in Section~\ref{sec:results}. A well-formed $p$-value made the artefact look
more credible, not less.

The sixth concerns the green-test gate itself. Our Python implementation retained the
passing tests and discarded the failing ones, as T2 specifies; our Java implementation
rejected the entire function if any single test failed. The two languages were measuring
different constructs under one name, and the Java side contradicted the definition we had
registered in advance. We corrected it and report both settings throughout
(Section~\ref{sec:results}); the correction is discussed under conclusion validity below,
because it moves results in our own favour.

None of these produce an error. Each returns a plausible number. We therefore report the four
nested tiers of Section~\ref{sec:tiers} rather than any single flag, and treat only T4 ---
fault detection within the target's line range --- as evidence that a suite tests its target.

\paragraph{Mutation operators.}
Our operator set is deliberately small and identical across both languages, so that the two
are measured with the same instrument. It is weaker than a production mutation framework, so
absolute mutation scores are not comparable with studies using PIT or \texttt{mutmut}; only
within-study comparisons are meaningful. The per-function mutant cap differs between the two
language pipelines (20 for Java, 30 for Python), an inconsistency inherited from separate
implementations. It binds two Java functions --- CJ-004 and CJ-026, both sitting exactly at
the 20-mutant ceiling under every Java condition --- and no Python function, the largest
there being CP-025 at 24 candidate sites against a cap of 30. Since T4 asks only whether at
least one mutant is killed, the cap bounds the resolution of the continuous score for those
two functions rather than their tier outcome.

\subsection{Internal Validity}

\paragraph{Sampling.}
Our own miner mis-attributed one Java target: \texttt{LocalTime.getField} is declared
\texttt{protected}, but the parser matched a same-named declaration in a nested class and
recorded it as public, so it passed the public-only filter. An audit against the real
declaration line found this to be the only such case in 60 Java functions; we exclude it from
the analysis. We record it here because it is precisely the class-attribution error this study
is about, reproduced in our own tooling --- knowing about a failure mode does not confer
immunity to it.

\paragraph{Concurrency.}
Mutation analysis overwrites source files. Running it directly on a shared source tree
corrupts any concurrent measurement reading that tree, silently. We mutate an isolated copy
placed ahead of the original on the import path, and verify after each run that the original
tree is unmodified.

\subsection{External Validity}

\paragraph{Source exhaustion in Python.}
After all testability criteria, only 64 Python functions in the two projects qualify, of which
34 already appeared in our earlier dataset. A strictly held-out Python replication is therefore
limited to roughly 30 functions without adding new repositories. Our Python sample is also
drawn from two web-oriented libraries and may not generalise to other domains.

\paragraph{Complexity distribution.}
Although the band is $5 \le CC \le 10$, the sample concentrates at the lower end (56 of 120
functions have $CC = 5$). Conclusions about the upper half of the band rest on fewer
observations.

\paragraph{Model and configuration.}
A single model at a single temperature was used. We do not claim the findings transfer to
other models, and the deterministic setting means we do not estimate run-to-run variance.

\subsection{Conclusion Validity}

\paragraph{Baseline configuration.}
Three limitations constrain how far the baseline comparison can be pushed. EvoSuite was run
with a classpath containing only the module's own compiled classes: with the full dependency
classpath its bytecode reader aborts in a way that produces no output while reporting success,
and we were unable to identify the specific offending artifact. EvoSuite may therefore be
weaker here than in a fully configured environment. Seven Java targets in one multi-module
project could not be recompiled to Java~11 bytecode and fall outside EvoSuite's coverage.
Finally, our Randoop version differs from the one used in our earlier study, so Randoop numbers
are compared only within this study.

\paragraph{Pre-registration and two deviations.}
Research questions, hypotheses, tool versions, budgets and the two-branch Pynguin design were
committed before any baseline was run.

The first deviation added a condition. On inspection, the enriched prompt we had planned to
use turned out to lack the worked exemplar, making it zero-shot rather than one-shot and
differing from the baseline in two respects at once. We added a corrected condition that
holds the exemplar fixed and varies only the target specification, recorded the change in the
pre-registration with its date, and left the original hypothesis unaltered. The uncorrected
variant is reported only as a secondary observation.

The second deviation changed the instrument after results had been seen, and it changed it in
a direction favourable to us, so we state it plainly. The Java green-test gate described
above was corrected to match the registered definition of T2. Six functions under the
baseline prompt and eight under the target-specified prompt were affected, against two for
EvoSuite and one for Randoop --- so the correction helps the model more than it helps its
competitors. Three safeguards apply. The pre-correction results were written to the
repository and committed \emph{before} the corrected measurement was run, so the timestamps
preclude selecting between them afterwards. Both settings are reported side by side for every
Java comparison. And we state no conclusion that does not hold under both; where the two
disagree, we report the comparison as inconclusive. The Python measurements were not re-run,
having matched the registered definition from the start.
